%% Part II, school Report #2. Grading: 10% on-going, 5% defence. Budget: 3 pages.
%% Draft by Trung (Refs #2). Tables hold checked facts; Trung writes the prose around them.
\chapter{Project Management Plan}\label{ch:management}

\section{Team Work}\label{sec:teamwork}


\subsection{Team Structure and Roles}\label{sec:roles}
\guide{Describe the structure of the team and the roles of each member: specific tasks, responsibilities, and communication methods between members.}
\Cref{tab:roles} lists the role and the main areas of work of each member.

\begin{table}[htb]
  \centering
  \caption{Roles and main areas of work of the four members}\label{tab:roles}
  \begin{tabularx}{\textwidth}{@{}p{0.22\textwidth} L L@{}}
    \toprule
    {Member} & {Role} & {Main areas of work}\\
    \midrule
    Đoàn Duy Long & Team leader & Market-data server (opentrading-mcp), benchmark (opentrading-benchmark), analysis skills\\
    Trần Huy Tuấn & \TBC{confirm that the FinDaHuman account is Tuấn} & Harness core, terminal interface, releases\\
    Đoàn Minh Hiếu & Member & Memory of the harness; later tools, permissions and the agent loop\\
    Hoàng Hiếu Trung & Member & MCP client, web search and fetch tools, stop gate, compaction, evaluation thresholds; Parts I to III of this report\\
    \bottomrule
  \end{tabularx}
  \fignote{Source: authors of merged pull requests in the three repositories; thesis assignment in the group's tracking issue.}
\end{table}

\subsection{Communication Plan}\label{sec:communication}
\guide{Describe how and how often project information will be communicated among team members and advisor (and other stakeholders).}

\section{Project Management Approach}\label{sec:approach}


\subsection{WBS}\label{sec:wbs}
\guide{Provide detailed descriptions for each task or duty within the WBS, including objectives, required resources, and assigned responsibilities.}
\Cref{tab:wbs} breaks the project into five work packages, and \cref{tab:milestones} places their milestones in time.

\begin{table}[htb]
  \centering
  \caption{Work packages, their outputs and their leads}\label{tab:wbs}
  \begin{tabularx}{\textwidth}{@{}p{0.2\textwidth} L p{0.22\textwidth} p{0.14\textwidth}@{}}
    \toprule
    {Package} & {Main tasks} & {Output} & {Lead}\\
    \midrule
    Agent harness & Agent loop, system prompt, tools, permissions, sub-agents, ultradesk, terminal interface & OpenTrading version 0.5.1 & \TBC{Tuấn}\\
    Memory & Memory store, investor constraints, recall & Memory of OpenTrading & Hiếu\\
    Market-data server & Market-data tools, shared answer contract, cloud deployment & opentrading-mcp & Long\\
    Evaluation & Task set, frozen data snapshot, comparison of three harnesses & opentrading-benchmark & Long\\
    Report & Parts I to IX, figures, references & This report & All members\\
    \bottomrule
  \end{tabularx}
  \fignote{Source: the group's time plan and tracking issue.}
\end{table}

%% Evidence: E36, E37.
\begin{table}[htb]
  \centering
  \caption{Milestones of the project}\label{tab:milestones}
  \begin{tabularx}{\textwidth}{@{}p{0.08\textwidth} p{0.22\textwidth} L@{}}
    \toprule
    {Week} & {Dates (2026)} & {Milestone}\\
    \midrule
    0 & 31 Aug -- 6 Sep & Topic chosen: multi-agent systems for financial analysis\\
    2 & 14 Sep -- 20 Sep & Survey of Claude Code, OpenCode, Codex and Cursor; keywords fixed\\
    4 & 28 Sep -- 4 Oct & Demo I: first data server, first skills, basic memory\\
    7 & 19 Oct -- 25 Oct & Demo II\\
    8 & 23 Oct & Version 0.2.1, the first public release\\
    10 & 9 Nov -- 15 Nov & Benchmark study or16 completed\\
    12 & 24 Nov, 25 Nov & Versions 0.5.0 and 0.5.1\\
    12 & 28 Nov & Review 1.1 documents submitted\\
    13 & 30 Nov -- 3 Dec & Review 1.1\\
    \bottomrule
  \end{tabularx}
  \fignote{Source: the group's time plan and minutes; release page of OpenTrading. Week dates are inferred from week 13 being Review 1.1. \TBC{check the week dates against the FAP schedule}}
\end{table}

\subsection{Risk Management}\label{sec:risk}
\guide{Describe potential risks during the project execution and preventive measures or response plans for each risk.}
\Cref{tab:risks} lists risks that occurred during the project and the response to each.

%% Evidence: E38--E42.
\begin{table}[htb]
  \centering
  \caption{Risks that occurred and the response to each}\label{tab:risks}
  \begin{tabularx}{\textwidth}{@{}p{0.18\textwidth} L L@{}}
    \toprule
    {Risk} & {What happened} & {Response}\\
    \midrule
    Release failure & Version 0.2.0 was never published: the package registry refused its name & Republished as version 0.2.1 under a scoped package name\\
    Data library expiry & A release candidate of the vnstock data library was switched off on 11 November 2026 & Moved to version 6.0.0, valid until 10 November 2027\\
    Provider quota & Free model tiers allow as few as 20 requests per day & Limits documented for users; a paid key planned as fallback\\
    Infrastructure fault & 132 benchmark trials of OpenTrading failed before calling the model, after a broken link during disk clean-up & Trials rerun and the rerun recorded in the study report\\
    Credential expiry & The publishing token expires by 21 January 2027 & Move to trusted publishing before then\\
    \bottomrule
  \end{tabularx}
  \fignote{Source: release notes and documentation of OpenTrading; documentation of opentrading-mcp; report of study or16.}
\end{table}

\subsection{Quality Management}\label{sec:quality}
\guide{Details on the measures and processes applied to ensure the quality of the project, including testing and validation of the product or project outcomes.}
\Cref{tab:quality-checks} lists the automated checks that every change passes before it is merged.

%% Evidence: E11, E13, E43.
\begin{table}[htb]
  \centering
  \caption{Automated checks run on each change}\label{tab:quality-checks}
  \begin{tabularx}{\textwidth}{@{}p{0.22\textwidth} p{0.26\textwidth} L@{}}
    \toprule
    {Check} & {When it runs} & {What it verifies}\\
    \midrule
    Type check and tests & Every pull request to OpenTrading, on Ubuntu and Windows & \num{2253} tests pass, none fail\\
    Shell tests & Every pull request, on macOS and Ubuntu & Arithmetic commands give the same result on each shell\\
    Build check & Every pull request & The program builds and passes its self-check\\
    Secret scan & Every pull request & No key or token appears in the history\\
    Installer checks & Releases and installer changes & Install, update and uninstall on three operating systems\\
    Data-server tests & Every pull request to opentrading-mcp & \num{1338} tests pass; \qty{71.69}{\percent} line coverage\\
    Benchmark tests & Every pull request to opentrading-benchmark & \num{1199} tests of the current study pass\\
    \bottomrule
  \end{tabularx}
  \fignote{Source: automated check runs of the three repositories.}
\end{table}

\subsection{Budget and Funding (Optional)}\label{sec:budget}
\guide{If applicable, describe the project budget and funding sources, with estimated costs and financial management throughout the project.}
\Cref{tab:costs} lists the recorded costs; the project had no external funding.

%% Evidence: E07, E44, E45.
\begin{table}[htb]
  \centering
  \caption{Recorded and planned costs}\label{tab:costs}
  \begin{tabularx}{\textwidth}{@{}L S[table-format=1.3] L@{}}
    \toprule
    {Item} & {Cost (USD)} & {Note}\\
    \midrule
    Model use of OpenTrading in study or16 & 1.498 & Lower bound: one trial has no recorded cost\\
    Model use of OpenCode in study or16 & 0.972 & \\
    Model use of DeepSeek Harness in study or16 & 0.944 & \\
    Hosting of the data server & 0.000 & Free plan of the cloud host\\
    \bottomrule
  \end{tabularx}
  \fignote{Source: report of study or16; deployment file of opentrading-mcp. A paid key of 10 to 15 USD was planned in week 8. \TBC{reconcile the 8 USD in the week-10 minutes with the 3.414 USD above}}
\end{table}

\subsection{Change Management Process}\label{sec:change}
\guide{Describe the process followed for any changes in the project scope, schedule, resources, or other requirements.}

\subsection{Closure and Evaluation}\label{sec:closure}
\guide{Describe the process for closing and evaluating the project upon completion, including how experiences and lessons learned will be summarized.}
